\documentclass[11pt]{article}

\usepackage[margin=1in]{geometry}
\usepackage[T1]{fontenc}
\usepackage{lmodern}
\usepackage{microtype}
\usepackage{amsmath,amssymb}
\usepackage{booktabs}
\usepackage{enumitem}
\usepackage{xcolor}
\usepackage{hyperref}

\definecolor{boleroblue}{HTML}{245A7A}
\definecolor{lightblue}{HTML}{EEF6FA}
\hypersetup{
  colorlinks=true,
  linkcolor=boleroblue,
  urlcolor=boleroblue,
  pdftitle={Conditioned LoRA in Bolero},
  pdfauthor={Bolero project}
}

\newcommand{\code}[1]{\texttt{\small #1}}
\newcommand{\Wbase}{W_{0}}
\newcommand{\condvec}{e}

\title{Conditioned LoRA in Bolero}
\author{Technical model note}
\date{July 2026}

\begin{document}

\maketitle

\begin{abstract}
Bolero is a cell-state-conditioned sequence-to-function model. It combines a frozen
Borzoi/Flashzoi DNA backbone with low-rank adapters whose weights are generated from a
cell-state embedding. The resulting effective network is different for each conditioning
vector while sharing the expensive sequence model across cells. This note describes the
conditioning path, the conditional LoRA parameterization, the paired signal-model training
recipe, and inference by collapsing the adapters into a fixed network.
\end{abstract}

\section{Model overview}

The model takes a long DNA sequence and a representation of the target cell state. The
default sequence input is one-hot DNA of length 524,288 bp. The output is a chromatin or
transcript track at 32 bp resolution. The backbone produces 16,384 internal output bins;
the default training crop is 16,352 bins (see \code{crop\_to\_length} in
\code{src/bolero/tl/model/borzoi/model.py}).

The main implementation is \code{BorzoiLoRA} in
\code{src/bolero/tl/model/borzoi/model\_lora.py}. The base checkpoint is loaded first,
and selected convolutional and linear layers are then replaced by LoRA modules.

\section{Conditional LoRA}

\subsection{From ordinary LoRA to conditional LoRA}

For a linear layer with frozen weight $\Wbase\in\mathbb{R}^{o\times i}$, ordinary LoRA
uses a learned, input-independent low-rank update. Bolero instead generates both low-rank
factors from the conditioning vector:

\begin{align}
A(\condvec) &\in \mathbb{R}^{r\times i},\\
B(\condvec) &\in \mathbb{R}^{o\times r},\\
W(\condvec) &= \Wbase + \frac{\alpha}{r}B(\condvec)A(\condvec).
\end{align}

Each factor is produced by a small MLP, implemented as an \code{EmbeddingMLP}. The MLPs
are shared across all examples; only their outputs vary with $\condvec$. This amortizes
the cost of learning a separate effective network for every cell state.

The implementation stores the first factor with shape $(r,i)$ and the second with shape
$(o,r)$, contracts them over the rank dimension, and adds the result to the frozen base
weight. The scale is $\alpha/r$; with the default configuration, $\alpha$ is obtained from
the rank and the configured LoRA scale.

\subsection{Explicit convolutional parameterization}

For a one-dimensional convolution with base kernel
\[
K_0 \in \mathbb{R}^{o\times(i/g)\times k},
\]
where $i$ and $o$ are the input and output channel counts, $g$ is the number of groups,
and $k$ is the kernel width, the implementation flattens the channel--kernel dimensions
before applying the low-rank factorization. For conditioning vector $\condvec$, its two
MLPs generate
\begin{align}
A_c(\condvec) &\in \mathbb{R}^{(rk)\times((i/g)k)},\\
B_c(\condvec) &\in \mathbb{R}^{o\times(rk)}.
\end{align}
The generated convolution update is
\[
\Delta K(\condvec) = \frac{\alpha}{r}\,
\operatorname{reshape}_{o,\,i/g,\,k}\!\left[
B_c(\condvec)A_c(\condvec)\right],
\]
and the effective kernel for one example is
\[
K(\condvec) = K_0 + \Delta K(\condvec).
\]
Thus, the factorization rank in the flattened representation is $rk$; the code's LoRA
rank parameter is $r$, and the kernel width is folded into both generated factors. For a
two-dimensional convolution, the analogous generator shapes are
\[
A_c(\condvec)\in\mathbb{R}^{(rk)\times((i/g)k)},\qquad
B_c(\condvec)\in\mathbb{R}^{(ok)\times(rk)},
\]
followed by reshaping $B_cA_c$ to $(o,i/g,k,k)$.

For a batch of $b$ examples, the one-dimensional implementation generates
$K(\condvec_1),\ldots,K(\condvec_b)$, concatenates them into a grouped-convolution weight
tensor
\[
K_{\mathrm{batch}}\in\mathbb{R}^{(bo)\times(i/g)\times k},
\]
reshapes the input $X\in\mathbb{R}^{b\times i\times L}$ to
$X'\in\mathbb{R}^{1\times(bi)\times L}$, and calls convolution with $bg$ groups. The
output is reshaped back to $Y\in\mathbb{R}^{b\times o\times L'}$. Each example therefore
uses its own conditional kernel while the batch is processed in one grouped operation.

The conditional layer implementation is in
\code{src/bolero/tl/generic/module\_lora\_cond.py}. Its two defining properties are:

\begin{enumerate}[leftmargin=2em]
  \item the pretrained base weight is frozen;
  \item the $A$ and $B$ generators are trainable and initialized so the initial LoRA
        update is zero (the final layer of the $B$ generator is zero-initialized).
\end{enumerate}

Zero initialization makes the converted model initially behave like the pretrained
backbone, before the adapters learn a cell-state-dependent correction.

\subsection{Which layers are adapted?}

The default preset is \code{all\_conditional}. Its principal ranks are:

\begin{center}
\begin{tabular}{ll}
\toprule
Backbone region & LoRA rank \\
\midrule
DNA convolution & 2 \\
Residual tower and first U-Net block & 6 \\
Transformer-associated layers & 30 \\
Horizontal and upsampling convolutions & 30 \\
Separable convolutions & 10 \\
Final joined convolutions & 60 \\
Output heads & 1 \\
\bottomrule
\end{tabular}
\end{center}

Most selected layers use conditional LoRA. The transformer query, key, value, relative-key,
and output projections are explicitly excluded from conditional LoRA; they receive ordinary
input-independent LoRA adapters instead. This exception is specified by
\code{exclude\_cond\_lora\_patterns} in
\code{src/bolero/tl/model/borzoi/model\_lora\_config.py}.

\section{Building the conditioning vector}

\subsection{What is the cell-state embedding?}

In the current embedding tutorial, the cell-state embedding is \emph{not} a set of
expression principal components. It is a learned latent representation from a joint
RNA--ATAC \textsc{MultiVI} model (implemented by \code{scvi-tools}). The worked tutorial
uses a 30-dimensional latent vector. MultiVI is trained on raw RNA counts for the top
5,000 highly variable genes and ATAC peak counts, with donor/experimental batch modeled as
a covariate. The trained model is then applied to all cells, and the latent representation
is stored as \code{X\_multivi} / \code{multivi.latent\_embedding.feather}.

SEACells uses this joint latent space to form homogeneous metacells. Bolero does not
usually feed individual noisy cells directly to the sequence model: the pseudobulk records
carry the embedding associated with each metacell or pseudobulk. A pseudobulk embedding is
typically the mean of the embeddings of its selected metacells. Therefore, the embedding
is a learned multimodal cell-state coordinate, not an expression-only PCA score. The
embedding dimensionality and model can be changed by the data-preparation workflow; the
LoRA input width is read from the resulting records.
The LoRA code itself is agnostic to the embedding's provenance: an expression-PCA vector
could be supplied through the same \code{emb\_key} interface, but that is not the embedding
used by the current MULTIVI-based tutorial and atlas recipe.

\subsection{Single-dataset training}

The paired pseudobulk pipeline samples a background pseudobulk $p_0$ and a target pseudobulk
$p_1$. Their cell-state embeddings are concatenated:

\[
\condvec_{\mathrm{cell}} = [\condvec_{p_0};\condvec_{p_1}].
\]

For example, two 30-dimensional metacell embeddings give a 60-dimensional input to the
LoRA generators. The background embedding describes the reference state, while the target
embedding describes the state whose track is to be predicted. The construction is performed
by \code{EnsemblePairedPseudobulker} in
\code{src/bolero/tl/pseudobulk/paired\_pseudobulk.py}.

Optional condition variables, such as tissue, developmental state, or TF activity, are
encoded separately. \code{ConditionEmbeddingModule} applies small encoders to the cell and
condition inputs and concatenates their outputs. If no condition variables are present, the
module is replaced by a no-effect module and the cell embedding is passed through unchanged.

\subsection{Multi-dataset training}

Different datasets can have different embedding spaces and condition schemas. The
multi-dataset model therefore uses:

\begin{itemize}[leftmargin=2em]
  \item a dataset-specific cell encoder;
  \item a dataset-specific condition encoder when conditions exist;
  \item a shared encoder for genome/species and other shared metadata.
\end{itemize}

The resulting dataset-specific and shared representations are concatenated into one vector
of fixed width before being supplied to all conditional-LoRA generators. This allows one
shared backbone to serve multiple datasets while retaining dataset-specific conditioning.

\section{Signal-model forward pass}

With \code{signal\_model=True}, the model receives the background signal from $p_0$ in
addition to DNA. The signal is transformed with $\log(1+s)$ and passed through a small
convolutional encoder. Its output is injected into the DNA representation at 32 bp
resolution. The subsequent Borzoi residual tower, U-Net path, transformer, and upsampling
path all receive the same conditioning vector and therefore use state-specific adapters.

The target output is produced by a 1D output head. For the default count head, a softplus
activation makes predictions nonnegative and training uses the Poisson-multinomial loss.
In paired training, the model can also predict the difference between target and background
signals with an auxiliary delta head and mean-squared-error loss.

At training time, the background signal is randomly omitted with probability
\code{nosignal\_prob} (default 0.2). This prevents the model from depending entirely on the
background track and encourages prediction from DNA plus conditioning. At inference time,
the signal input may be omitted; the DNA and conditioning paths remain available.

\section{Training and trainable parameters}

The trainer first loads the pretrained Borzoi checkpoint and freezes the backbone. It then
converts the configured layers to LoRA modules. The base convolution and linear weights
remain frozen, while the following components are trainable in the standard recipe:

\begin{itemize}[leftmargin=2em]
  \item conditional-LoRA generator MLPs;
  \item ordinary LoRA parameters in the attention projections;
  \item the cell/condition encoders;
  \item the output heads;
  \item the signal encoder and enabled auxiliary heads.
\end{itemize}

The optimizer is built from parameters whose \code{requires\_grad} flag is true. The
backbone is thus reused as a strong sequence prior, while the comparatively small adapter
and conditioning components learn cell-state-specific regulation.

\section{Collapsing for inference}

For a fixed conditioning vector $\condvec^*$, every conditional layer can be evaluated once:

\[
W^* = \Wbase + \frac{\alpha}{r}B(\condvec^*)A(\condvec^*).
\]

\code{BorzoiLoRA.collapse\_lora(embedding)} replaces the conditional layers with ordinary
PyTorch layers containing these fixed weights. The result is a conventional DNA-to-track
model for one cell state. It is useful for fast repeated prediction and for base-level DNA
attribution. A different collapsed model is needed for each conditioning vector whose
effective network is to be inspected.

\section{Implementation map}

\begin{center}
\begin{tabular}{p{0.38\linewidth}p{0.52\linewidth}}
\toprule
Component & Source file \\
\midrule
Conditional linear and convolutional LoRA & {\scriptsize\code{src/bolero/tl/generic/module\_lora\_cond.py}} \\
Condition encoders and MLP generators & {\scriptsize\code{src/bolero/tl/generic/module\_embedding.py}} \\
BorzoiLoRA construction and forward paths & {\scriptsize\code{src/bolero/tl/model/borzoi/model\_lora.py}} \\
Layer/rank presets & {\scriptsize\texttt{model\_lora\_config.py}} \\
Paired training forward pass & {\scriptsize\code{src/bolero/tl/model/borzoi/train.py}} \\
Pseudobulk pairing and embeddings & {\scriptsize\code{src/bolero/tl/pseudobulk/paired\_pseudobulk.py}} \\
\bottomrule
\end{tabular}
\end{center}

\end{document}
